\scanpage{55}
\mspage{54}

% slanted
\margin{\uline{25 Juin}}

Je voudrais revenir sur la notion d'\uline{atelier ensembliste}. Ce sera un atelier \unsure{soit} spacieux, soit polyédral, prenons le premier cas.

\uline{Proposition} Soit $\unsure{\mathcal A|\mathscr A|A}=(\unsure{\mathcal F|\mathscr F|F},\leq,\ll)$ un atelier, considérons le (co)morphisme canonique

\[
\unsure{\varphi|\phi}:\mathcal F\longrightarrow\mathrm{Fig}(\unsure{\mathcal L|\mathscr L})\qquad F\longmapsto\unsure{\mathrm{Malomb}|\mathrm{malomb}}(F)
\]

\inserted{Considérons les} Conditions : \deleted{\unsure{équivalentes}}

(i) L'hom. $\varphi$ est un \uline{plongement} d'ateliers (cf \unsure{CF|GF} \unsure{VI|II} p. 51) et l'image de $\mathcal F$ par $\varphi$ est un sous-atelier \uline{spacieux} (loc cit \inserted{p. 40} : si $F,G\in\varphi(\mathcal F)$ et $F\mathrel{>\mkern-13mu<}G$ \deleted{\unsure{entre|est}} dans $\mathrm{Fig}(\mathcal L)$, alors $F\mathrel{>\mkern-13mu<}G$ dans $\varphi(\mathcal F)$ i.e $F\cup G\in\deleted{\unsure{F}}\varphi(\mathcal F)$)

En d'autres termes :

% slanted
\margin{$\varphi$ plongement}

\[
\left\{\begin{aligned}
&\text{a) }\forall F\in\mathcal F,\quad F'\unsure{\mathrel{\vdash\mkern-6mu\dashv}|\longmapsto}\varphi(F')\ \text{induit une bijection}\\
&\hspace{6em}\mathrm{Ss\,fig}(F)\longrightarrow\mathrm{Ss\,fig}(\varphi(F))\\
&\text{b) }\varphi(F)\ll\varphi(G)\Rightarrow F\ll G\quad(\text{donc }\varphi\ \text{\unsure{injective|injection}})
\end{aligned}\right.
\]

% slanted
\margin{$\varphi$ plongement ``spacieux''}

\[
\text{c) Si }\varphi(F)\mathrel{>\mkern-13mu<}\varphi(G),\Rightarrow F\mathrel{>\mkern-13mu<}G
\]

(ii) $\mathcal F$ est \uline{spacieux} (au sens absolu) i.e. satisfait \deleted{\unsure{les}} l'un des quatre systèmes \unsure{équivalents|équivalent} de \unsure{conditions|condition} de la p. 45 \inserted{pour $\unsure{\Lambda|\wedge}=\mathcal L$} (p.ex. $\mathrm{At}\,L\,1$, $\mathrm{At}\,L\,2$, $\mathrm{At}\,L\,3$, $\mathrm{At}\,CL(\unsure{\mathrm{spac}|\mathrm{spa}})$), \uuline{et} pour $x,y\in\mathcal L$, $x\ne y$ \deleted{\unsure{on a}}\inserted{$\Rightarrow$} $x\deleted{\unsure{=}}\parallel y$. (donc $\mathrm{At}\,L\,3$ devient automatique)

Alors $(i)\Rightarrow(ii)$ et $(ii)\Rightarrow(i\,a)\deleted{\unsure{(b)|(c)}}$ (mais pas néc $(i\,b)$, \inserted{ni $(i,c)$,} ni l'injectivité de $\varphi$)

\uline{Dém} $(i)\Rightarrow(ii)$. \deleted{Comme} Commençons par $x,y\in\mathcal L$, $x\ne y\Rightarrow\deleted{\unsure{\varphi|y}}x\parallel y$ ou encore, \deleted{$(x,y)\in\mathcal L$}\inserted{$x,y\in\mathcal L$}$\Rightarrow x\mathrel{>\mkern-13mu<}y$. Comme $\varphi(x),\varphi(y)$ sont des \deleted{\unsure{familles}} figures ens. ponctuelles on a $\varphi(x)\mathrel{\deleted{\unsure{\times}}>\mkern-13mu<}\varphi(y)$ donc $x\mathrel{>\mkern-13mu<}y$ par c). $\mathrm{At}\,L\,1$ : Soit $X\in\unsure{M|\mathcal M}$, \deleted{soit $x$} comme $X\ne\varnothing$ on a $\varphi(\unsure{x|X})\ne\varnothing$, donc par définition de $\varphi$ \unsure{Cela|Celà} signifie $\mathrm{omb}(X)^{\circ}\ne\varnothing$, \deleted{\unsure{et|ce}} OK.

